% ch04g —— Wentworth《Higgs Bundles and Local Systems on Riemann Surfaces》
% §4.4.3 后段（引理 4.4.10 至引理 4.4.21）与 §4.4.4 前段
% 底本：工程文件/work/en/ch04.txt 第 2635–3204 行；印刷页 207–213（PDF 217–223）
% 本段编号公式：(4.55)–(4.64)

\begin{lemma}
\[
H^1\bigl(X,(\mathcal{L}\otimes\mathcal{K}^{n-j})\otimes\mathcal{R}_i^*\bigr)=
\begin{cases}
0, & \text{若 } i\ge j+1,\\
H^1(X,\mathcal{K}), & \text{若 } i=j.
\end{cases}
\]
\end{lemma}

\begin{proof}
固定 $j$，对 $i$ 作归纳。若 $i=n-1$，则 $\mathcal{R}_{n-1}=\mathcal{L}$，从而
\[
H^1\bigl(X,(\mathcal{L}\otimes\mathcal{K}^{n-j})\otimes\mathcal{R}_{n-1}^*\bigr)
=H^1(X,\mathcal{K}^{n-j})=
\begin{cases}
0, & \text{若 } n-j-1>0,\\
H^1(X,\mathcal{K}), & \text{若 } n=j+1.
\end{cases}
\]
现在，正合列 $0\to\mathcal{R}_i^*\to\mathcal{R}_{i-1}^*\to\mathcal{L}^*\otimes\mathcal{K}^{i-n}\to0$ 给出
\[
H^1\bigl(X,(\mathcal{L}\otimes\mathcal{K}^{n-j})\otimes\mathcal{R}_i^*\bigr)
\longrightarrow
H^1\bigl(X,(\mathcal{L}\otimes\mathcal{K}^{n-j})\otimes\mathcal{R}_{i-1}^*\bigr)
\longrightarrow
H^1(X,\mathcal{K}^{i-j})\longrightarrow0.
\]
由归纳假设，第一项消失，而后两项同构。引理得证。
\end{proof}

\begin{lemma}
\[
H^1\bigl(X,\mathcal{V}_j\otimes\mathcal{R}_i^*\bigr)=
\begin{cases}
0, & \text{若 } i\ge j+1,\\
H^1(X,\mathcal{K}), & \text{若 } i=j.
\end{cases}
\]
\end{lemma}

\begin{proof}
固定 $i$，对 $j$ 作归纳。此时 $\mathcal{V}_1\simeq\mathcal{L}\otimes\mathcal{K}^{n-1}$，故这一情形下的结论由引理 4.4.10 推出。其次，考虑正合列
\[
H^1(X,\mathcal{V}_{j-1}\otimes\mathcal{R}_i^*)\longrightarrow
H^1(X,\mathcal{V}_j\otimes\mathcal{R}_i^*)\longrightarrow
H^1(X,\mathcal{L}\otimes\mathcal{K}^{n-j}\otimes\mathcal{R}_i^*)\longrightarrow0.
\]
由归纳假设，第一项消失，于是第二、三项同构。再次由引理 4.4.10 即得结论。
\end{proof}

\begin{corollary}
\[
H^1\bigl(X,\mathcal{V}_{j-1}\otimes(\mathcal{L}\otimes\mathcal{K}^{n-j})^*\bigr)=H^1(X,\mathcal{K}).
\]
\end{corollary}

\begin{proof}
考虑正合列
\[
H^1(X,\mathcal{V}_{j-1}\otimes\mathcal{R}_j^*)\longrightarrow
H^1(X,\mathcal{V}_{j-1}\otimes\mathcal{R}_{j-1}^*)\longrightarrow
H^1\bigl(X,\mathcal{V}_{j-1}\otimes(\mathcal{L}_{j-1}\otimes\mathcal{K}^{n-j})^*\bigr)\longrightarrow0.
\]
由引理 4.4.11，第一项消失，而第二项同构于 $H^1(X,\mathcal{K})$。
\end{proof}

\begin{annote}
底本在此处把商线丛印作 $\mathcal{L}_{j-1}\otimes\mathcal{K}^{n-j}$，但同段其他公式以及引理 4.4.13 均作 $\mathcal{L}\otimes\mathcal{K}^{n-j}$（即 $\mathcal{V}_j/\mathcal{V}_{j-1}$）；此处照排原文，疑为笔误，见疑点 Q097。
\end{annote}

\begin{lemma}
扩张 $0\to\mathcal{V}_{j-1}\to\mathcal{V}_j\to\mathcal{L}\otimes\mathcal{K}^{n-j}\to0$ 是非分裂的。
\end{lemma}

\begin{proof}
考虑交换图
\begin{equation}
\begin{tikzcd}[row sep=2.4em, column sep=3.4em]
& H^1\bigl(X,\mathcal{V}_{j-1}\otimes\mathcal{R}_j^*\bigr)\arrow[d,"g"']\\
H^0\bigl(X,\mathcal{R}_{j-1}\otimes\mathcal{R}_{j-1}^*\bigr)\arrow[r,"{I\mapsto\lbrack\beta\rbrack}"]\arrow[d]
& H^1\bigl(X,\mathcal{V}_{j-1}\otimes\mathcal{R}_{j-1}^*\bigr)\arrow[d]\\
H^0\bigl(X,\mathcal{R}_{j-1}\otimes(\mathcal{L}\otimes\mathcal{K}^{n-j})^*\bigr)\arrow[r]
& H^1\bigl(X,\mathcal{V}_{j-1}\otimes(\mathcal{L}\otimes\mathcal{K}^{n-j})^*\bigr)\arrow[d]\\
H^0\bigl(X,(\mathcal{L}\otimes\mathcal{K}^{n-j})\otimes(\mathcal{L}\otimes\mathcal{K}^{n-j})^*\bigr)\arrow[u]\arrow[ur,"{I\mapsto\lbrack\alpha\rbrack}"']
& 0
\end{tikzcd}
\tag{4.55}
\end{equation}
由 (4.5) 之后的评注，$[\alpha]$ 是扩张 $0\to\mathcal{V}_{j-1}\to\mathcal{V}_j\to\mathcal{L}\otimes\mathcal{K}^{n-j}\to0$ 的扩张类，而 $[\beta]$ 是扩张 $0\to\mathcal{V}_{j-1}\to\mathcal{V}\to\mathcal{R}_{j-1}\to0$ 的扩张类。由引理 4.4.11，$g$ 是单射。循上边缘的定义追踪可知 $[\alpha]=g[\beta]$。最后，由于 $\mathcal{V}$ 带有全纯连接，且由引理 4.4.9 有 $\deg\mathcal{V}_{j-1}\ne0$，故由定理 4.3.5 得 $[\beta]\ne0$。
\end{proof}

\begin{annote}
底本把 $[\beta]$ 所属的扩张写作 $0\to\mathcal{V}_{j-1}\to\mathcal{V}\to\mathcal{R}_{j-1}\to0$，但 $\mathcal{V}$ 的秩为 $n$、$\mathcal{R}_{j-1}$ 的秩为 $1$，秩不匹配；按上下文（$[\beta]\in H^1(X,\mathcal{V}_{j-1}\otimes\mathcal{R}_{j-1}^*)$）中间项应为 $\mathcal{V}_j$。此处照排原文，见疑点 Q098。
\end{annote}

最后，我们可以叙述关于底全纯结构唯一性的结果。

\begin{proposition}
设 $(\mathcal{V},\nabla)$ 为 $\mathrm{SL}_n$-oper。则 $\mathcal{V}$ 上的 oper 结构由 $\mathcal{L}=\mathcal{V}/\mathcal{V}_{n-1}$ 唯一确定。特别地，丛 $\mathcal{V}$ 的同构类在 $\operatorname{Op}_n$ 的每个连通分支上固定。
\end{proposition}

\begin{proof}
由引理 4.4.9，$\mathcal{V}_1=\mathcal{L}\otimes\mathcal{K}^{n-1}$，故它被确定。由推论 4.4.12 与引理 4.4.13，每个 $\mathcal{V}_j$ 都作为引理 4.4.13 中那个序列的唯一非分裂扩张，被 $\mathcal{V}_{j-1}$ 逐步确定。如此继续直到 $j=n$，即得第一个论断。第二个论断同样成立，因为由引理 4.4.9 还有 $\mathcal{L}^n\simeq\mathcal{K}^{-n(n-1)/2}$，从而可能的 $\mathcal{L}$ 的集合是离散的。
\end{proof}

\begin{corollary}
把 oper 映为其单值表示的映射给出一个嵌入 $\operatorname{Op}_n\to\mathfrak{M}_B^{(n)}$。
\end{corollary}

\begin{proof}
固定表示 $\rho\colon\pi\to\mathrm{SL}_n(\C)$，并设（在共轭意义下）$\rho$ 同时是 oper $(\mathcal{V}_\rho,\nabla_1)$ 与 $(\mathcal{V}_\rho,\nabla_2)$ 的单值表示。鉴于命题 4.4.14，只需证明线丛 $\mathcal{L}$ 被 $\rho$ 唯一确定。设 $\mathcal{L}$ 与 $\mathcal{M}$ 是次数为 $-(n-1)(g-1)$ 的线丛，使得 $H^0(X,\mathcal{V}_\rho^*\otimes\mathcal{L})\ne\{0\}$ 且 $H^0(X,\mathcal{V}_\rho^*\otimes\mathcal{M})\ne\{0\}$。设 $\{\mathcal{V}_i\}$ 是 $(\mathcal{V}_\rho,\nabla_1)$ 的 oper 结构，并设 $\mathcal{V}_\rho/\mathcal{V}_{n-1}=\mathcal{L}$。若 $\mathcal{L}$ 与 $\mathcal{M}$ 不同构，则由
\[
0\to\mathcal{L}^*\otimes\mathcal{M}\to\mathcal{V}_\rho^*\otimes\mathcal{M}\to\mathcal{V}_{n-1}^*\otimes\mathcal{M}\to0
\]
可得 $H^0(X,\mathcal{V}_{n-1}^*\otimes\mathcal{M})\ne\{0\}$。现在，对 $j\le n-1$ 有 $\deg\mathcal{R}_j^*\otimes\mathcal{M}<0$，故反复应用这一论证，便推出 $H^0(X,\mathcal{V}_1^*\otimes\mathcal{M})\ne\{0\}$。但 $\mathcal{V}_1^*\otimes\mathcal{M}=\mathcal{L}^*\otimes\mathcal{M}\otimes\mathcal{K}^{1-n}$ 的次数也为负，于是得到矛盾。
\end{proof}

\begin{remark}
命题 4.4.14 中的线丛 $\mathcal{L}$ 恰有 $n^{2g}$ 种可能。这些选择给出 $\operatorname{Op}_n$ 各连通分支的标记。如同注 4.4.2，它们恰好对应于把 $\mathrm{PSL}_n(\C)$ 中的单值表示提升为 $\mathrm{SL}_n(\C)$ 中的表示的 $n^{2g}$ 种方式。为简单起见，今后我们总取 $\mathcal{L}=\mathcal{K}^{-(n-1)/2}$，其中当 $n$ 为偶数时假定已固定一个 $\mathcal{K}^{1/2}$ 的选取。
\end{remark}

\paragraph{oper 与微分方程}

我们首先说明如何从一个实现在 $\mathcal{K}^{1-q}$ 中的局部系统得到 oper，此处 $n=2q-1$。于是，设已给正合列 (4.52)，并令 $\mathcal{V}=\mathcal{V}_n=\calO\otimes_\C\mathbf{V}$。对 $k=1,\dots,n-1$，定义
\[
\mathcal{V}_{n-k}=\Bigl\{\sum_{i=1}^n f_i\otimes\mathbf{v}_i:\sum_{i=1}^n f_i^{(j)}\varphi(\mathbf{v}_i)=0,\ j=0,\dots,k-1\Bigr\}.
\]
则 $\mathcal{V}_{n-k}\subset\mathcal{V}$ 是凝聚子层（coherent subsheaf），且有正合列
\begin{equation}
0\to\mathcal{V}_{n-k-1}\to\mathcal{V}_{n-k}\to\mathcal{K}^{1-q+k}\to0.
\tag{4.56}
\end{equation}
定义 4.4.7 中的性质 (i) 显然满足。此外，鉴于 (4.56)，连接 $\nabla$ 诱导 $\calO$-线性映射 $\mathcal{V}_{n-k-1}\to\mathcal{V}_{n-k}/\mathcal{V}_{n-k-1}\otimes\mathcal{K}\simeq\mathcal{K}^{2-q+k}$，它由
\[
\sum_{i=1}^n f_i\otimes\mathbf{v}_i\mapsto\sum_{i=1}^n f_i^{(k+1)}\varphi(\mathbf{v}_i)
\]
给出，且这是层的同构。故性质 (ii) 也成立。

反之，设 $\mathcal{V}$ 是秩 $n$ 的全纯丛，带有全纯连接 $\nabla$，且容许一个 oper 结构。由引理 4.4.9，有 $\mathcal{V}/\mathcal{V}_{n-1}\simeq\mathcal{K}^{1-q}$。由此推出，对任一 $\mathrm{SL}_n$-oper（我们仍假定 $\mathcal{L}=\mathcal{K}^{-(n-1)/2}$），有 $H^0(X,\mathcal{V}^*\otimes\mathcal{K}^{1-q})\ne\{0\}$。由于 oper 的单值表示由命题 4.4.8 是不可约的，定理 4.4.5 的假设满足，故 $(\mathcal{V},\nabla)$ 实现在 $\mathcal{K}^{1-q}$ 中。

\begin{theorem}[{Beilinson--Drinfeld, [3]}]
上述嵌入给出 $\operatorname{Op}_n$ 的连通分支与（仿射）Hitchin 底空间（Hitchin base）$\bigoplus_{j=2}^n H^0(X,\mathcal{K}^j)$ 之间的同构。
\end{theorem}

\begin{corollary}[{Teleman, [63]}]
微分方程 (4.52)（或 (4.53)）的单值表示永远不是酉的。
\end{corollary}

\begin{proof}
若 $\rho$ 是单值表示，则由上述对应与引理 4.4.9 可知，$\mathcal{V}_\rho$ 是不稳定丛。但于是由定理 4.2.6 的容易方向（见命题 4.2.16），$\mathcal{V}_\rho$ 不能容许平坦酉连接。
\end{proof}

\paragraph{oper 与模空间}

本小节的主要目标是证明

\begin{theorem}
映射 $\operatorname{Op}_n\hookrightarrow\mathfrak{M}_B^{(n)}$ 是紧合嵌入（proper embedding）。
\end{theorem}

由 Harder--Narasimhan 型（Harder--Narasimhan type）的上半连续性，本定理是下列命题的直接推论（见第 4.2.1 小节）。

\begin{proposition}
在带有全纯连接的丛当中，oper 具有严格极大的 Harder--Narasimhan 型。
\end{proposition}

我们从下述引理开始。

\begin{lemma}
带有 oper 结构的丛 $\mathcal{V}$ 的 Harder--Narasimhan 滤过由 oper 滤过本身给出。
\end{lemma}

\begin{proof}
只需证明，对每个 $j=0,\dots,n-1$，$\mathcal{V}_{j+1}/\mathcal{V}_j$ 是 $\mathcal{V}/\mathcal{V}_j$ 的极大失稳子层（maximally destabilizing subsheaf）。为此，记 $\mu_{\max}(\mathcal{V}/\mathcal{V}_j)$ 为子层 $\mathcal{F}\subset\mathcal{V}/\mathcal{V}_j$（$0<\operatorname{rank}\mathcal{F}<\operatorname{rank}(\mathcal{V}/\mathcal{V}_j)$）的最大斜率。我们作归纳假设
\[
\mu_{\max}(\mathcal{V}/\mathcal{V}_j)=\mu(\mathcal{V}_{j+1}/\mathcal{V}_j)=(n-1)(g-1)-j(2g-2).
\]
注意当 $j=n-1$ 时它平凡地成立。现在设 $j\le n-2$，并设 $\mathcal{F}\subset\mathcal{V}/\mathcal{V}_j$ 是极大失稳子层。则 $\mathcal{F}$ 是半稳定的，并且由正合列
\[
0\to\mathcal{V}_{j+1}/\mathcal{V}_j\to\mathcal{V}/\mathcal{V}_j\to\mathcal{V}/\mathcal{V}_{j+1}\to0
\]
与归纳假设，有
\[
\mu(\mathcal{F})\ge\mu(\mathcal{V}_{j+1}/\mathcal{V}_j)>\mu(\mathcal{V}_{j+2}/\mathcal{V}_{j+1})=\mu_{\max}(\mathcal{V}/\mathcal{V}_{j+1}).
\]
由此推出诱导映射 $\mathcal{F}\to\mathcal{V}/\mathcal{V}_{j+1}$ 必为零。因此 $\mathcal{F}\simeq\mathcal{V}_{j+1}/\mathcal{V}_j$，归纳假设对 $j$ 成立。证毕。
\end{proof}

\begin{proof}[命题 4.4.20 的证明]
（参看 [40, 定理 5.3.1]。）设 $(\mathcal{V},\nabla)$ 是带全纯连接的不稳定丛。我断言只需假定 $\nabla$ 不可约。事实上，秩为 $1$ 的情形无需证明。设结论对小于 $n$ 的秩已证，并设 $(\mathcal{V},\nabla)$ 可约。由于 Harder--Narasimhan 型上半连续，可假定存在分裂 $(\mathcal{V},\nabla)=(\mathcal{V}_1,\nabla_1)\oplus(\mathcal{V}_2,\nabla_2)$，其中 $n_i=\operatorname{rank}\mathcal{V}_i\ge1$。于是由归纳假设，只需假定 $\mathcal{V}_i$ 具有秩 $n_i$ 的 oper 的 Harder--Narasimhan 型。的确，若非如此，我们就能在不改变 $\mathcal{V}$ 的斜率次序的前提下改变 $\mathcal{V}_i$ 的 Harder--Narasimhan 型，使 $\mathcal{V}$ 具有更大的 Harder--Narasimhan 型。令
\begin{equation}
\mu_i=\mu_i^{(n)}=\mu(\mathcal{K}^{q-i})=(n+1-2i)(g-1)
\tag{4.57}
\end{equation}
为秩 $n$ 的 oper 的 Harder--Narasimhan 型（见引理 4.4.21 与 4.4.9）。若 $\lambda_i$ 是斜率 $\{\mu_i^{(n_1)},\mu_j^{(n_2)}\}$ 的一个重排，则需证明
\begin{equation}
\sum_{i=1}^k\lambda_i\le\sum_{i=1}^k\mu_i^{(n)},
\tag{4.58}
\end{equation}
对一切 $k=1,\dots,n$ 成立，且对某个 $k$ 取严格不等式。设 $n_1\ge n_2$。在不改变斜率次序的前提下，我们可以依次从排在前面的项 $\mu_i^{(n)}$、$\lambda_i=\mu_i^{(n_1)}$ 中减去偶数，只要 $2i\le n_1-n_2$，再把这些整数加到末尾的项上，只要 $n_1+n_2+2\le2i$。注意如此得到的 $\{\mu_i\}$ 与 $\{\lambda_i\}$ 中首、末斜率的重数相等，在求和中相消。这就把问题化归为 $n_1=n_2$ 或 $n_1=n_2+1$ 的情形，此时容易验证 (4.58)。

至此可以假定 $(\mathcal{V},\nabla)$ 不可约。oper 的 Harder--Narasimhan 型由 (4.57) 给出。设 $\mathcal{V}_{i-1}\subset\mathcal{V}_i$（$i=1,\dots,\ell$）是 $\mathcal{V}$ 的 Harder--Narasimhan 滤过，并令 $\lambda_i=\mu(\mathcal{V}_i/\mathcal{V}_{i-1})$。令 $n_i=\operatorname{rank}(\mathcal{V}_i/\mathcal{V}_{i-1})$、$d_i=n_i\lambda_i$。则只需证明
\begin{equation}
\sum_{i=1}^j n_i\lambda_i\le\sum_{i=1}^{\operatorname{rank}(\mathcal{V}_j)}\mu_i,
\tag{4.59}
\end{equation}
对 $j=1,\dots,\ell$ 成立。左边恰是 $\deg\mathcal{V}_j$，而右边是
\[
\sum_{i=1}^{\operatorname{rank}(\mathcal{V}_j)}(n+1-2i)(g-1)=(g-1)\operatorname{rank}(\mathcal{V}_j)\bigl(n-\operatorname{rank}(\mathcal{V}_j)\bigr).
\]
于是 (4.59) 等价于
\begin{equation}
\deg\mathcal{V}_j\le(g-1)\Bigl(\sum_{i=1}^j n_i\Bigr)\Bigl(n-\sum_{i=1}^j n_i\Bigr).
\tag{4.60}
\end{equation}
反复应用命题 4.3.4，可得 $\lambda_j\le\lambda_{j+1}+2g-2$、$\lambda_j\le\lambda_{j+2}+2(2g-2)$、$\lambda_j\le\lambda_{j+i}+i(2g-2)$ 以及 $\lambda_j\le\lambda_\ell+(\ell-j)(2g-2)$，其中 $i\le\ell-j$ 任意。这意味着
\[
\frac{n_{j+1}}{n_j}d_j\le d_{j+1}+(2g-2)n_{j+1},
\]
\[
\frac{n_{j+i}}{n_j}d_j\le d_{j+i}+i(2g-2)n_{j+i},
\]
\[
\frac{n_\ell}{n_j}d_j\le d_\ell+(\ell-j)(2g-2)n_{j+1},
\]
\begin{annote}
底本此三行不等式的第三行末尾印作 $n_{j+1}$，但按 $\lambda_j\le\lambda_\ell+(\ell-j)(2g-2)$ 乘以 $n_\ell$ 的推导应为 $n_\ell$；此处照排原文，见疑点 Q099。
\end{annote}
由此得到
\begin{equation}
\Bigl(\sum_{i=1}^{\ell-j}n_{i+j}\Bigr)\frac{d_j}{n_j}\le\sum_{i=1}^{\ell-j}d_{i+j}+(2g-2)\sum_{i=1}^{\ell-j}in_{i+j}.
\tag{4.61}
\end{equation}
先考虑 $j=1$ 的情形。此时 (4.61) 成为
\[
\Bigl(\sum_{i=2}^{\ell}n_i\Bigr)\frac{d_1}{n_1}\le\sum_{i=2}^{\ell}d_i+(2g-2)\sum_{i=2}^{\ell}(i-1)n_i,
\]
\[
(n-n_1)\frac{d_1}{n_1}\le-d_1+(2g-2)\sum_{i=2}^{\ell}(i-1)n_i,
\]
从而
\begin{equation}
d_1\le\frac{n_1}{n}(2g-2)\sum_{i=2}^{\ell}(i-1)n_i.
\tag{4.62}
\end{equation}
我们断言
\begin{equation}
\frac{2}{n}\sum_{i=2}^{\ell}(i-1)n_i\le n-n_1=\sum_{i=2}^{\ell}n_i.
\tag{4.63}
\end{equation}
注意它与 (4.62) 合起来便证明了 $j=1$ 时的 (4.60)。为证此断言，令 $r_i=n_i-1\ge0$。则 (4.63) 成为
\[
2\sum_{i=2}^{\ell}(i-1)(r_i+1)\le n\sum_{i=2}^{\ell}(r_i+1),
\]
\[
2\sum_{i=2}^{\ell}(i-1)r_i+\ell(\ell-1)\le\Bigl[\sum_{i=2}^{\ell}(r_i+1)\Bigr]^2+n_1\sum_{i=2}^{\ell}(r_i+1),
\]
只要下式成立即可
\[
2\sum_{i=2}^{\ell}(i-1)r_i+\ell(\ell-1)\le\Bigl[\sum_{i=2}^{\ell}r_i+(\ell-1)\Bigr]^2+(\ell-1),
\]
而后式在消去两边同类项之后，又只要
\[
\sum_{i=2}^{\ell}(i-1)r_i\le\sum_{i=2}^{\ell}(\ell-1)r_i
\]
成立即可。后一不等式因 $r_i\ge0$ 显然成立。故 (4.63) 成立。

我们现在作归纳。设 (4.60) 对 $j$ 成立，我们证明它对 $j+1$ 也成立。把 $j$ 的 (4.60) 与 $j+1$ 的 (4.61) 相加，得
\[
\begin{aligned}
\deg\mathcal{V}_{j+1}&=\deg\mathcal{V}_j+d_{j+1}\\
&\le(g-1)\sum_{i=1}^j n_i\sum_{i=1}^{\ell-j}n_{i+j}-\frac{n_{j+1}}{\sum_{i=2}^{\ell-j}n_{i+j}}\sum_{i=1}^{j+1}d_i+\frac{n_{j+1}}{\sum_{i=2}^{\ell-j}n_{i+j}}(2g-2)\sum_{i=1}^{\ell-j-1}in_{i+j+1},
\end{aligned}
\]
\[
\frac{\sum_{i=1}^{\ell-j}n_{i+j}}{\sum_{i=2}^{\ell-j}n_{i+j}}\deg\mathcal{V}_{j+1}\le(g-1)\sum_{i=1}^j n_i\sum_{i=1}^{\ell-j}n_{i+j}+\frac{n_{j+1}}{\sum_{i=2}^{\ell-j}n_{i+j}}(2g-2)\sum_{i=1}^{\ell-j-1}in_{i+j+1},
\]
\[
\frac{\deg\mathcal{V}_{j+1}}{n-\sum_{i=1}^{j+1}n_i}\le(g-1)\sum_{i=1}^j n_i+\frac{n_{j+1}}{\bigl(\sum_{i=1}^{\ell-j}n_{i+j}\bigr)\bigl(\sum_{i=2}^{\ell-j}n_{i+j}\bigr)}(2g-2)\sum_{i=1}^{\ell-j-1}in_{i+j+1},
\]
其中从第一个不等式过渡到第二个不等式时，用到了 $\deg\mathcal{V}_{j+1}=\sum_{i=1}^{j+1}d_i$。于是只需证明
\[
2\sum_{i=1}^{\ell-j-1}in_{i+j+1}\le\sum_{i=1}^{\ell-j}n_{i+j}\sum_{i=2}^{\ell-j}n_{i+j}.
\]
用上面定义的 $r_i$ 来表示，它成为
\[
2\sum_{i=2}^{\ell-j}(i-1)(r_{i+j}+1)\le(r_{j+1}+1)\sum_{i=2}^{\ell-j}(r_{i+j}+1)+\Bigl(\sum_{i=2}^{\ell-j}(r_{i+j}+1)\Bigr)^2,
\]
\[
2\sum_{i=2}^{\ell-j}(i-1)r_{i+j}+(\ell-j)(\ell-j-1)\le(r_{j+1}+1)\sum_{i=2}^{\ell-j}r_{i+j}+(r_{j+1}+1)(\ell-j-1)+\Bigl(\sum_{i=2}^{\ell-j}r_{i+j}+(\ell-j-1)\Bigr)^2.
\]
但这可由
\[
2\sum_{i=2}^{\ell-j}(i-1)r_{i+j}\le2\sum_{i=2}^{\ell-j}(\ell-j-1)r_{i+j}
\]
推出，而后者显然成立。这就完成了 Harder--Narasimhan 型极大性的证明。下面我们证明：若 $(\mathcal{V},\nabla)$ 的 Harder--Narasimhan 型极大，则滤过 $\{\mathcal{V}_i\}$ 是 oper 结构。事实上，考虑 $\calO$-线性映射 $\nabla\colon\mathcal{V}_i\to\mathcal{V}/\mathcal{V}_{i+1}\otimes\mathcal{K}$。由注 4.2.9，$\mathcal{V}_i$ 的商的极小斜率是 $\mu_i=\mu(\mathcal{V}_i/\mathcal{V}_{i-1})$，而 $\mathcal{V}/\mathcal{V}_{i+1}\otimes\mathcal{K}$ 的子层的极大斜率为
\[
\mu(\mathcal{V}_{i+2}/\mathcal{V}_{i+1}\otimes\mathcal{K})=\mu_{i+2}+2g-2=\mu_{i+1}=\mu_i-(2g-2)<\mu_i.
\]
因此上述映射必须为零，且 $\nabla\mathcal{V}_i\subset\mathcal{V}_{i+1}\otimes\mathcal{K}$。由连接的不可约性，$\mathcal{V}_i/\mathcal{V}_{i-1}\to\mathcal{V}_{i+1}/\mathcal{V}_i\otimes\mathcal{K}$ 非零。由于这些都是同次数的线丛，该映射是同构。因此定义 4.4.7 中的条件 (i) 与 (ii) 满足。证毕。
\end{proof}

\subsection{Eichler--Shimura 同构}

让我们更详细地回到例 4.4.6。对 $q\in\frac{1}{2}\Z$，记 $V_q$ 为 $\mathrm{SL}_2(\C)$ 的（$2q-1$ 维）不可约表示。设 $\rho\colon\pi\to\mathrm{SL}_2(\C)$ 是 $X$ 上某个射影连接的单值表示（的提升）。对自旋结构（Spin structure）的某个选取，我们可以把局部系统 $\mathbf{V}_\rho$ 实现在 $\mathcal{K}^{-1/2}$ 中。对 $q\ge\frac{3}{2}$，记 $\mathbf{V}_q$ 为把 $\rho$ 与表示 $V_q$ 复合所得的局部系统，
\[
\rho^{(n)}\colon\pi\to\mathrm{SL}_2(\C)\to\mathrm{SL}(V_q).
\]
则 $\mathbf{V}_q$ 实现在 $\mathcal{K}^{1-q}$ 中，并且有
\begin{equation}
0\to\mathbf{V}_q\to\mathcal{K}^{1-q}\xrightarrow{\ D\ }\mathcal{K}^q\to0.
\tag{4.64}
\end{equation}
